\paragraph{Disclosed cash constrains a history without identifying it.}
Importing a verified native book begins with supplied account and transaction
facts. A public cash-flow statement supplies a different kind of evidence.
The retained Apple annual statement reports opening cash of USD~30,737 million,
a net decrease of USD~794 million and closing cash of USD~29,943 million.
Those three figures constrain a possible movement history, but they do not
identify its gross receipts, payments, transaction dates or legal accounts.
The source record therefore remains distinct from a modeller's proposed
explanation of the movement.

Consider two deliberately hypothetical borrowing and repayment histories.
Let \(O\) and \(C\) denote opening and closing cash, and let \(B\) and \(P\)
denote cash borrowing and principal repayment within the annual period, all
in USD millions. Assume that these are the only cash movements in the example.
The resulting cash relation is
\[
  C=O+B-P,\qquad
  29{,}943=30{,}737+0-794
           =30{,}737+100-894.
\]
The second history has an additional borrowing of 100 and an additional
repayment of 100. It reaches the same closing cash while preserving a different
gross event sequence. Both examples start with an assumed debt of 1,000 and
balancing capital; the event dates and finite amount bounds are also supplied
assumptions. These are counterexamples to identification from the selected
cash figures, not disclosures of Apple's borrowing or histories shown to fit
its other statements. The selected source model additionally assumes that
unrounded values satisfy the cash bridge and uses the earlier proposed
nearest-million rounding rule. That rule is not silently attributed to the issuer.

The optional \texttt{RetainedSourceHistoryAttachment/1} makes this distinction
operational. Its request supplies the complete retained statement request,
a typed history for each selected cash bridge, and explicit mappings from
source quantities to history fields. It recomputes the source report, retains
all 256 observations and 143 relation dispositions, and constructs each
attachment from the original source identities, dates, monetary scale and
basis context. The existing history and measurement methods then check both
complete constraint systems and every supplied mapping. A history that fixes
both annual cash movements to zero is feasible by itself but fails when
attached to the annual decrease. Missing rounding or mapping premises remain
unresolved; an exhausted computation budget yields an unknown result rather
than an infeasibility conclusion.

The interim bridge illustrates a second distinction. Its opening 29,943,
increase 356 and closing 30,299 permit hypothetical borrowing and repayment
pairs \((356,0)\) and \((456,100)\). The annual closing and interim opening
share a date and displayed amount, but their retained basis contexts differ.
The two attachments consequently remain separate: equality of the displayed
boundary cannot authorize a common cross-period history. Every candidate also
retains the unresolved actual-transaction and unselected-disclosure obligations.
A selected-cash witness therefore supplies no admitted legal book, cash-flow
classification or compatibility conclusion for the whole filing.

Eleven complete exchanges and their semantic replay check this composition,
including alternative exact witnesses, changed sources, missing mappings,
incompatible contexts and unresolved premises. Fifteen focused tests and
65 existing source/history tests pass. The implementation and full exchanges
are retained under the following path, split for legibility:
\begin{quote}\small\ttfamily
data/accounting-master/\\
retained-source-history-20261025/
\end{quote}
The next step for a possible company history is to retain these source
conditions when selecting its opening and continuation.
